\documentclass[11pt]{article}

\usepackage[utf8]{inputenc}
\usepackage{fullpage}
\usepackage{amsmath, amssymb, bm, cite, epsfig, psfrag}
\usepackage{graphicx}
\usepackage{float}
\usepackage{amsthm}
\usepackage{amsfonts}
\usepackage{listings}
\usepackage{cite}
\usepackage{hyperref}
\usepackage{tikz}
\usepackage{enumerate}
\usepackage{listings}
\usepackage{mathtools}
\lstloadlanguages{Python}
\usetikzlibrary{shapes,arrows}
%\usetikzlibrary{dsp,chains}

\DeclareFixedFont{\ttb}{T1}{txtt}{bx}{n}{9} % for bold
\DeclareFixedFont{\ttm}{T1}{txtt}{m}{n}{9}  % for normal
% Defining colors
\usepackage{color}
\definecolor{deepblue}{rgb}{0,0,0.5}
\definecolor{deepred}{rgb}{0.6,0,0}
\definecolor{deepgreen}{rgb}{0,0.5,0}
\definecolor{backcolour}{rgb}{0.95,0.95,0.92}

%\restylefloat{figure}
%\theoremstyle{plain}      \newtheorem{theorem}{Theorem}
%\theoremstyle{definition} \newtheorem{definition}{Definition}

\def\del{\partial}
\def\ds{\displaystyle}
\def\ts{\textstyle}
\def\beq{\begin{equation}}
\def\eeq{\end{equation}}
\def\beqa{\begin{eqnarray}}
\def\eeqa{\end{eqnarray}}
\def\beqan{\begin{eqnarray*}}
\def\eeqan{\end{eqnarray*}}
\def\nn{\nonumber}
\def\binomial{\mathop{\mathrm{binomial}}}
\def\half{{\ts\frac{1}{2}}}
\def\Half{{\frac{1}{2}}}
\def\N{{\mathbb{N}}}
\def\Z{{\mathbb{Z}}}
\def\Q{{\mathbb{Q}}}
\def\R{{\mathbb{R}}}
\def\C{{\mathbb{C}}}
\def\argmin{\mathop{\mathrm{arg\,min}}}
\def\argmax{\mathop{\mathrm{arg\,max}}}
%\def\span{\mathop{\mathrm{span}}}
\def\diag{\mathop{\mathrm{diag}}}
\def\x{\times}
\def\limn{\lim_{n \rightarrow \infty}}
\def\liminfn{\liminf_{n \rightarrow \infty}}
\def\limsupn{\limsup_{n \rightarrow \infty}}
\def\GV{Guo and Verd{\'u}}
\def\MID{\,|\,}
\def\MIDD{\,;\,}

\newtheorem{proposition}{Proposition}
\newtheorem{definition}{Definition}
\newtheorem{theorem}{Theorem}
\newtheorem{lemma}{Lemma}
\newtheorem{corollary}{Corollary}
\newtheorem{assumption}{Assumption}
\newtheorem{claim}{Claim}
\def\qed{\mbox{} \hfill $\Box$}
\setlength{\unitlength}{1mm}

\def\bhat{\widehat{b}}
\def\ehat{\widehat{e}}
\def\phat{\widehat{p}}
\def\qhat{\widehat{q}}
\def\rhat{\widehat{r}}
\def\shat{\widehat{s}}
\def\uhat{\widehat{u}}
\def\ubar{\overline{u}}
\def\vhat{\widehat{v}}
\def\xhat{\widehat{x}}
\def\xbar{\overline{x}}
\def\zhat{\widehat{z}}
\def\zbar{\overline{z}}
\def\la{\leftarrow}
\def\ra{\rightarrow}
\def\MSE{\mbox{\small \sffamily MSE}}
\def\SNR{\mbox{\small \sffamily SNR}}
\def\SINR{\mbox{\small \sffamily SINR}}
\def\arr{\rightarrow}
\def\Exp{\mathbb{E}}
\def\var{\mbox{var}}
\def\Tr{\mbox{Tr}}
\def\tm1{t\! - \! 1}
\def\tp1{t\! + \! 1}
\def\Tm1{T\! - \! 1}
\def\Tp1{T\! + \! 1}


\def\Xset{{\cal X}}

\newcommand{\one}{\mathbf{1}}
\newcommand{\abf}{\mathbf{a}}
\newcommand{\bbf}{\mathbf{b}}
\newcommand{\dbf}{\mathbf{d}}
\newcommand{\ebf}{\mathbf{e}}
\newcommand{\gbf}{\mathbf{g}}
\newcommand{\hbf}{\mathbf{h}}
\newcommand{\pbf}{\mathbf{p}}
\newcommand{\pbfhat}{\widehat{\mathbf{p}}}
\newcommand{\qbf}{\mathbf{q}}
\newcommand{\qbfhat}{\widehat{\mathbf{q}}}
\newcommand{\rbf}{\mathbf{r}}
\newcommand{\rbfhat}{\widehat{\mathbf{r}}}
\newcommand{\sbf}{\mathbf{s}}
\newcommand{\sbfhat}{\widehat{\mathbf{s}}}
\newcommand{\ubf}{\mathbf{u}}
\newcommand{\ubfhat}{\widehat{\mathbf{u}}}
\newcommand{\utildebf}{\tilde{\mathbf{u}}}
\newcommand{\vbf}{\mathbf{v}}
\newcommand{\vbfhat}{\widehat{\mathbf{v}}}
\newcommand{\wbf}{\mathbf{w}}
\newcommand{\wbfhat}{\widehat{\mathbf{w}}}
\newcommand{\xbf}{\mathbf{x}}
\newcommand{\xbfhat}{\widehat{\mathbf{x}}}
\newcommand{\xbfbar}{\overline{\mathbf{x}}}
\newcommand{\ybf}{\mathbf{y}}
\newcommand{\zbf}{\mathbf{z}}
\newcommand{\zbfbar}{\overline{\mathbf{z}}}
\newcommand{\zbfhat}{\widehat{\mathbf{z}}}
\newcommand{\Ahat}{\widehat{A}}
\newcommand{\Abf}{\mathbf{A}}
\newcommand{\Bbf}{\mathbf{B}}
\newcommand{\Cbf}{\mathbf{C}}
\newcommand{\Bbfhat}{\widehat{\mathbf{B}}}
\newcommand{\Dbf}{\mathbf{D}}
\newcommand{\Gbf}{\mathbf{G}}
\newcommand{\Hbf}{\mathbf{H}}
\newcommand{\Ibf}{\mathbf{I}}
\newcommand{\Kbf}{\mathbf{K}}
\newcommand{\Pbf}{\mathbf{P}}
\newcommand{\Phat}{\widehat{P}}
\newcommand{\Qbf}{\mathbf{Q}}
\newcommand{\Rbf}{\mathbf{R}}
\newcommand{\Rhat}{\widehat{R}}
\newcommand{\Sbf}{\mathbf{S}}
\newcommand{\Ubf}{\mathbf{U}}
\newcommand{\Vbf}{\mathbf{V}}
\newcommand{\Wbf}{\mathbf{W}}
\newcommand{\Xhat}{\widehat{X}}
\newcommand{\Xbf}{\mathbf{X}}
\newcommand{\Ybf}{\mathbf{Y}}
\newcommand{\Zbf}{\mathbf{Z}}
\newcommand{\Zhat}{\widehat{Z}}
\newcommand{\Zbfhat}{\widehat{\mathbf{Z}}}
\def\alphabf{{\boldsymbol \alpha}}
\def\betabf{{\boldsymbol \beta}}
\def\betabfhat{{\widehat{\bm{\beta}}}}
\def\epsilonbf{{\boldsymbol \epsilon}}
\def\mubf{{\boldsymbol \mu}}
\def\lambdabf{{\boldsymbol \lambda}}
\def\etabf{{\boldsymbol \eta}}
\def\xibf{{\boldsymbol \xi}}
\def\taubf{{\boldsymbol \tau}}
\def\sigmahat{{\widehat{\sigma}}}
\def\thetabf{{\bm{\theta}}}
\def\thetabfhat{{\widehat{\bm{\theta}}}}
\def\thetahat{{\widehat{\theta}}}
\def\mubar{\overline{\mu}}
\def\muavg{\mu}
\def\sigbf{\bm{\sigma}}
\def\etal{\emph{et al.}}
\def\Ggothic{\mathfrak{G}}
\def\Pset{{\mathcal P}}
\newcommand{\bigCond}[2]{\bigl({#1} \!\bigm\vert\! {#2} \bigr)}
\newcommand{\BigCond}[2]{\Bigl({#1} \!\Bigm\vert\! {#2} \Bigr)}
\newcommand{\tran}{^{\text{\sf T}}}
\newcommand{\herm}{^{\text{\sf H}}}
\newcommand{\bkt}[1]{{\langle #1 \rangle}}
\def\Norm{{\mathcal N}}
\newcommand{\vmult}{.}
\newcommand{\vdiv}{./}


\def\hid{\textsc{\tiny H}}
\def\out{\textsc{\tiny O}}

% Python style for highlighting
\newcommand\pythonstyle{\lstset{
language=Python,
backgroundcolor=\color{backcolour},
commentstyle=\color{deepgreen},
basicstyle=\ttm,
otherkeywords={self},             % Add keywords here
keywordstyle=\ttb\color{deepblue},
emph={MyClass,__init__},          % Custom highlighting
emphstyle=\ttb\color{deepred},    % Custom highlighting style
stringstyle=\color{deepgreen},
%frame=tb,                         % Any extra options here
showstringspaces=false            %
}}

% Python environment
\lstnewenvironment{python}[1][]
{
\pythonstyle
\lstset{#1}
}
{}

% Python for external files
\newcommand\pythonexternal[2][]{{
\pythonstyle
\lstinputlisting[#1]{#2}}}

% Python for inline
\newcommand\pycode[1]{{\pythonstyle\lstinline!#1!}}

\begin{document}

\title{Introduction to Machine Learning\\
Neural Networks Problems}
\author{Prof. Sundeep Rangan and Yao Wang}
\date{}

\maketitle

% Submit answers to only problems 1--3.  But, make sure you know how to do all problems.

\begin{enumerate}



\item Consider a neural network
on a 3-dimensional input $\xbf=(x_1,x_2, x_3)$ with weights and biases:
\[
    W^\hid = \left[ \begin{array}{ccc} 1 & 0 & 1  \\ 0 & 1  & 1 \\ 1 & 1 & 0 \\ 1 & 1 & 1\\ \end{array} \right], \quad
    b^\hid = \left[ \begin{array}{c} 0 \\ 0 \\ -1 \\ 1 \end{array} \right] \quad
    W^\out = [1, 1, -1, -1], \quad b^\out = -1.5.
\]
Assume the network uses the
threshold activation function \eqref{eq:gact_ht}
\beq \label{eq:gact_ht}
    g_{\rm act}(z) = \begin{cases}
        1, & \mbox{if } z \geq 0 \\
        0, & \mbox{if } z < 0.
        \end{cases}
\eeq
 and the
threshold output function \eqref{eq:gout_thresh}:
\beq \label{eq:gout_thresh}
    \hat{y} =  \begin{cases}
        1 & z^\out \geq 0 \\
        0 & z^\out < 0.
        \end{cases}
\eeq



\begin{enumerate}[(a)]
\item Write the components of $\zbf^\hid$ and $\ubf^\hid$ as a function
of $(x_1,x_2, x_3)$.  For each component $j$, indicate where in the $(x_1,x_2, x_3)$
hyperplane $u^\hid_j=1$.

\textbf{Solution:}
We have $\zbf^\hid = W^\hid \xbf + \bbf^\hid$:
\begin{align*}
z_1^\hid &= x_1 + x_3 \\
z_2^\hid &= x_2 + x_3 \\
z_3^\hid &= x_1 + x_2 - 1 \\
z_4^\hid &= x_1 + x_2 + x_3 + 1
\end{align*}
Since $u_j^\hid = g_{\rm act}(z_j^\hid)$, we have:
\begin{align*}
u_1^\hid &= \begin{cases} 1 & \text{if } x_1 + x_3 \geq 0 \\ 0 & \text{otherwise} \end{cases} \\
u_2^\hid &= \begin{cases} 1 & \text{if } x_2 + x_3 \geq 0 \\ 0 & \text{otherwise} \end{cases} \\
u_3^\hid &= \begin{cases} 1 & \text{if } x_1 + x_2 \geq 1 \\ 0 & \text{otherwise} \end{cases} \\
u_4^\hid &= \begin{cases} 1 & \text{if } x_1 + x_2 + x_3 \geq -1 \\ 0 & \text{otherwise} \end{cases}
\end{align*}
The hyperplanes where $u_j^\hid = 1$ are:
\begin{itemize}
\item $u_1^\hid = 1$: $x_1 + x_3 \geq 0$
\item $u_2^\hid = 1$: $x_2 + x_3 \geq 0$
\item $u_3^\hid = 1$: $x_1 + x_2 \geq 1$
\item $u_4^\hid = 1$: $x_1 + x_2 + x_3 \geq -1$
\end{itemize}

\item Write $z^\out$ as a function of $(x_1,x_2, x_3)$.  In what region is
$\hat{y}=1$? (You can describe in mathematical formulae).

\textbf{Solution:}
We have $z^\out = W^\out \ubf^\hid + b^\out = u_1^\hid + u_2^\hid - u_3^\hid - u_4^\hid - 1.5$.

Since each $u_j^\hid \in \{0,1\}$, $z^\out$ can take values in $\{-3.5, -2.5, -1.5, -0.5, 0.5, 1.5, 2.5, 3.5\}$.

$\hat{y} = 1$ when $z^\out \geq 0$, which occurs when:
\[
u_1^\hid + u_2^\hid - u_3^\hid - u_4^\hid \geq 1.5
\]
Since the left side is an integer, this means:
\[
u_1^\hid + u_2^\hid - u_3^\hid - u_4^\hid \geq 2
\]
This is satisfied when:
\begin{itemize}
\item $u_1^\hid = u_2^\hid = u_4^\hid = 1$ and $u_3^\hid = 0$, i.e., $x_1 + x_3 \geq 0$, $x_2 + x_3 \geq 0$, $x_1 + x_2 < 1$, and $x_1 + x_2 + x_3 \geq -1$
\item $u_1^\hid = u_2^\hid = u_3^\hid = u_4^\hid = 1$, i.e., $x_1 + x_3 \geq 0$, $x_2 + x_3 \geq 0$, $x_1 + x_2 \geq 1$, and $x_1 + x_2 + x_3 \geq -1$
\end{itemize}
The feasible regions are:
\begin{align*}
\hat{y} = 1 \text{ when: } &(x_1 + x_3 \geq 0) \land (x_2 + x_3 \geq 0) \land (x_1 + x_2 < 1) \land (x_1 + x_2 + x_3 \geq -1) \\
&\text{OR } (x_1 + x_3 \geq 0) \land (x_2 + x_3 \geq 0) \land (x_1 + x_2 \geq 1) \land (x_1 + x_2 + x_3 \geq -1)
\end{align*}
\end{enumerate}



\item Consider a neural network used for regression with a scalar input $x$ and
scalar target $y$,
\begin{align*}
    & z_{j}^\hid = W^\hid_{j}x + b^\hid_j, \quad
    u_{j}^\hid = \max\{0, z^\hid_j\},
    \quad j=1,\ldots,N_h \\
    &
    z^\out = \sum_{k=1}^{N_h} W^\out_{k}u^\hid_{k} + b^\out,
    \quad
    \hat{y} = g_{\rm out}(z^\out).
\end{align*}
The hidden weights and biases are:
\[
    \Wbf^\hid = \left[ \begin{array}{c} -1  \\ 1  \\  1 \end{array} \right], \quad
    \bbf^\hid = \left[ \begin{array}{c} -1 \\ 1 \\ -2 \end{array} \right].
\]

\begin{enumerate}[(a)]

\item What is the number $N_h$ of hidden units?  For each $j=1,\ldots,N_h$,
draw $u^\hid_j$ as a function of $x$ over some suitable range of values $x$.
You may draw them on one plot, or on
multiple plots.

\textbf{Solution:}
$N_h = 3$ (from the dimension of $\Wbf^\hid$).

For each $j$:
\begin{align*}
u_1^\hid(x) &= \max\{0, -x - 1\} = \begin{cases} -x - 1 & \text{if } x \leq -1 \\ 0 & \text{if } x > -1 \end{cases} \\
u_2^\hid(x) &= \max\{0, x + 1\} = \begin{cases} 0 & \text{if } x < -1 \\ x + 1 & \text{if } x \geq -1 \end{cases} \\
u_3^\hid(x) &= \max\{0, x - 2\} = \begin{cases} 0 & \text{if } x < 2 \\ x - 2 & \text{if } x \geq 2 \end{cases}
\end{align*}
These are ReLU functions with breakpoints at $x = -1$, $x = -1$, and $x = 2$ respectively.

\item Since the network is for regression, you may choose
the activation function $g_{\rm out}(z^\out)$ to be linear.
Given training data $(x_i,y_i)$, $i=1,\ldots,N$, formally define the loss function you would use
to train the network?

\textbf{Solution:}
For regression with linear output activation, we use the mean squared error (MSE) loss:
\[
L = \frac{1}{N} \sum_{i=1}^N (y_i - \hat{y}_i)^2 = \frac{1}{N} \sum_{i=1}^N \left(y_i - \left(\sum_{k=1}^{N_h} W_k^\out u_k^\hid(x_i) + b^\out\right)\right)^2
\]

\item Using the output activation and loss function selected in part (b), set up the formulation to determine
output weights and bias, $\Wbf^\out$, $b^\out$,
for the training data below.  You should be able to find a closed form solution. Write a few line of python code to solve the problem.
\begin{center}
\begin{tabular}{|c|c|c|c|c|c|} \hline
$x_i$ & -2 & -1 & 0 & 3 & 3.5\\ \hline
$y_i$ & 0 & 0   & 1 & 3 & 3 \\ \hline
\end{tabular}
\end{center}

\textbf{Solution:}
First, compute $\ubf^\hid(x_i)$ for each training point:
\begin{align*}
\ubf^\hid(-2) &= [1, 0, 0]^T \\
\ubf^\hid(-1) &= [0, 0, 0]^T \\
\ubf^\hid(0) &= [0, 1, 0]^T \\
\ubf^\hid(3) &= [0, 4, 1]^T \\
\ubf^\hid(3.5) &= [0, 4.5, 1.5]^T
\end{align*}
Let $\Ubf = [\ubf^\hid(x_1), \ldots, \ubf^\hid(x_5)]^T$ (5×3 matrix) and $\ybf = [0, 0, 1, 3, 3]^T$.
We solve: $\ybf = \Ubf \Wbf^\out + b^\out \one$, where $\one$ is a vector of ones.
This is a linear system: $[\Ubf, \one] [\Wbf^\out; b^\out] = \ybf$.

\begin{python}
import numpy as np

# Training data
x = np.array([-2, -1, 0, 3, 3.5])
y = np.array([0, 0, 1, 3, 3])

# Compute hidden layer outputs
W_hid = np.array([[-1], [1], [1]])
b_hid = np.array([[-1], [1], [-2]])
u_hid = np.maximum(0, W_hid @ x.reshape(1, -1) + b_hid).T

# Set up linear system: [U, ones] @ [W_out; b_out] = y
U = np.hstack([u_hid, np.ones((5, 1))])
params = np.linalg.lstsq(U, y, rcond=None)[0]
W_out = params[:3]
b_out = params[3]

print(f"W_out = {W_out}")
print(f"b_out = {b_out}")
\end{python}

\item
 Based on your solution for the output weights and bias,
draw $\hat{y}$ vs.\ $x$ over some suitable range of values $x$.  Write a few line of python code to do this.

\textbf{Solution:}
\begin{python}
import numpy as np
import matplotlib.pyplot as plt

# Parameters (from previous solution)
W_hid = np.array([[-1], [1], [1]])
b_hid = np.array([[-1], [1], [-2]])
W_out = np.array([0, 0.5, 1])  # Example solution
b_out = 0.5  # Example solution

# Evaluate over range
x_plot = np.linspace(-3, 5, 1000)
u_hid_plot = np.maximum(0, W_hid @ x_plot.reshape(1, -1) + b_hid)
y_hat = W_out @ u_hid_plot + b_out

plt.plot(x_plot, y_hat)
plt.xlabel('x')
plt.ylabel('$\\hat{y}$')
plt.title('Network output vs input')
plt.grid(True)
plt.show()
\end{python}

\item Write a function \pycode{predict} to output $\hat{y}$ given a vector of inputs
\pycode{x}.  Assume \pycode{x} is a vector representing a batch of samples and
\pycode{yhat} is a vector with the corresponding outputs.  Use the activation function
you selected in part (b), but your function should take the weights and biases
for both layers as inputs.  Clearly state any assumptions on the formats for
the weights and biases.  Also, to receive full credit, you must not use any for loops.

\textbf{Solution:}
\textbf{Assumptions:}
\begin{itemize}
\item \pycode{x} is a 1D array of shape $(N,)$ where $N$ is the batch size
\item \pycode{W_hid} is a 1D array of shape $(N_h,)$
\item \pycode{b_hid} is a 1D array of shape $(N_h,)$
\item \pycode{W_out} is a 1D array of shape $(N_h,)$
\item \pycode{b_out} is a scalar
\end{itemize}

\begin{python}
import numpy as np

def predict(x, W_hid, b_hid, W_out, b_out):
    """
    Predict output for batch of inputs.
    
    Args:
        x: 1D array of shape (N,) - batch of inputs
        W_hid: 1D array of shape (N_h,) - hidden layer weights
        b_hid: 1D array of shape (N_h,) - hidden layer biases
        W_out: 1D array of shape (N_h,) - output layer weights
        b_out: scalar - output layer bias
    
    Returns:
        yhat: 1D array of shape (N,) - predictions
    """
    # Hidden layer: z_hid shape (N_h, N)
    z_hid = W_hid[:, np.newaxis] * x[np.newaxis, :] + b_hid[:, np.newaxis]
    u_hid = np.maximum(0, z_hid)  # ReLU activation, shape (N_h, N)
    
    # Output layer: linear
    yhat = W_out @ u_hid + b_out  # shape (N,)
    
    return yhat
\end{python}
\end{enumerate}

\pagebreak
\item \label{prob:cg_exp}
Consider a neural network that takes each input $\xbf$ and produces
a prediction $\hat{y}$ given
\begin{align} \label{eq:nnmax}
\begin{split}
    z_j  &= \sum_{k=1}^{N_i}W_{jk}x_k + b_j, \quad u_j = 1/(1+\exp(-z_j)),
    \quad j = 1,\ldots,M,\\
    \hat{y} &= \frac{ \sum_{j=1}^{M} a_j u_j}{ \sum_{j=1}^M u_j},
\end{split}
\end{align}
where $M$ is the number of hidden units and is fixed (i.e.\ not trainable).
To train the model, we get training data $(\xbf_i,y_i)$, $i=1,\ldots,N$.
\begin{enumerate}[(a)]
\item Rewrite the equations \eqref{eq:nnmax} for the
batch of inputs $\xbf_i$ from the training data.  That is, correctly
add the indexing $i$, to the equations.

\textbf{Solution:}
\begin{align}
z_{ij} &= \sum_{k=1}^{N_i} W_{jk} x_{ik} + b_j, \quad u_{ij} = \frac{1}{1+\exp(-z_{ij})}, \quad j = 1,\ldots,M, \quad i = 1,\ldots,N \\
\hat{y}_i &= \frac{\sum_{j=1}^{M} a_j u_{ij}}{\sum_{j=1}^M u_{ij}}, \quad i = 1,\ldots,N
\end{align}

\item If we use a loss function,
\[
    L = \sum_{i=1}^N (y_i - \hat{y}_i)^2,
\]
draw the computation graph describing the mapping from $\xbf_i$ and parameters
to $L$.
Indicate which nodes are trainable parameters.

\textbf{Solution:}
The computation graph is:
\begin{align*}
\xbf_i &\rightarrow \zbf_i = \Wbf \xbf_i + \bbf \quad \text{(trainable: } \Wbf, \bbf \text{)} \\
\zbf_i &\rightarrow \ubf_i = \sigma(\zbf_i) \quad \text{(sigmoid)} \\
\ubf_i &\rightarrow \hat{y}_i = \frac{\sum_{j=1}^M a_j u_{ij}}{\sum_{j=1}^M u_{ij}} \quad \text{(weighted average)} \\
\hat{y}_i, y_i &\rightarrow L_i = (y_i - \hat{y}_i)^2 \\
L_i &\rightarrow L = \sum_{i=1}^N L_i
\end{align*}
Trainable parameters: $W_{jk}$, $b_j$ (and $a_j$ if trainable, but stated as fixed).

\item Compute the gradient  $\partial L/\partial \hat{y}_i$
for all $i$.

\textbf{Solution:}
\[
\frac{\partial L}{\partial \hat{y}_i} = \frac{\partial}{\partial \hat{y}_i} \sum_{k=1}^N (y_k - \hat{y}_k)^2 = -2(y_i - \hat{y}_i)
\]
So $\partial L/\partial \hat{\bf y} = -2(\ybf - \hat{\ybf})$ where $\hat{\ybf} = [\hat{y}_1, \ldots, \hat{y}_N]^T$.

\item Suppose that, in backpropagation, we have computed $\partial L/\partial \hat{y}_i$
for all $i$, represented as $\partial L/\partial \hat{\bf y}$. Describe how to compute the components of the gradient
$\partial L/\partial \ubf$.

\textbf{Solution:}
We have $\hat{y}_i = \frac{\sum_{j=1}^M a_j u_{ij}}{S_i}$ where $S_i = \sum_{j=1}^M u_{ij}$.

Using the quotient rule:
\[
\frac{\partial \hat{y}_i}{\partial u_{ij}} = \frac{a_j S_i - \sum_{k=1}^M a_k u_{ik}}{S_i^2} = \frac{a_j - \hat{y}_i}{S_i}
\]

By the chain rule:
\[
\frac{\partial L}{\partial u_{ij}} = \frac{\partial L}{\partial \hat{y}_i} \frac{\partial \hat{y}_i}{\partial u_{ij}} = \frac{\partial L}{\partial \hat{y}_i} \frac{a_j - \hat{y}_i}{S_i}
\]

In matrix form, if $\partial L/\partial \hat{\ybf}$ is a vector of length $N$ and $\Ubf$ is $N \times M$:
\[
\frac{\partial L}{\partial \Ubf} = \left(\frac{\partial L}{\partial \hat{\ybf}} \odot \frac{1}{\Sbf}\right) (\abf - \hat{\ybf})^T
\]
where $\odot$ is element-wise multiplication, $\Sbf = [S_1, \ldots, S_N]^T$, and the outer product gives the $N \times M$ gradient matrix.

\item
Suppose that we have computed the gradient $\partial L/\partial \ubf$, describe how would you compute the gradient $\partial L/\partial \zbf$.

\textbf{Solution:}
Since $u_{ij} = \sigma(z_{ij}) = 1/(1+\exp(-z_{ij}))$, we have:
\[
\frac{\partial u_{ij}}{\partial z_{ij}} = \sigma(z_{ij})(1 - \sigma(z_{ij})) = u_{ij}(1 - u_{ij})
\]

By the chain rule:
\[
\frac{\partial L}{\partial z_{ij}} = \frac{\partial L}{\partial u_{ij}} \frac{\partial u_{ij}}{\partial z_{ij}} = \frac{\partial L}{\partial u_{ij}} u_{ij}(1 - u_{ij})
\]

In matrix form:
\[
\frac{\partial L}{\partial \Zbf} = \frac{\partial L}{\partial \Ubf} \odot \Ubf \odot (1 - \Ubf)
\]
where all operations are element-wise.

\item
Suppose that we have computed the gradient $\partial L/\partial \zbf$, describe how would you compute the gradient
 $\partial L/\partial W_{jk}$ and $\partial L/\partial b_j$.

\textbf{Solution:}
Since $z_{ij} = \sum_{k=1}^{N_i} W_{jk} x_{ik} + b_j$, we have:
\[
\frac{\partial z_{ij}}{\partial W_{jk}} = x_{ik}, \quad \frac{\partial z_{ij}}{\partial b_j} = 1
\]

By the chain rule:
\[
\frac{\partial L}{\partial W_{jk}} = \sum_{i=1}^N \frac{\partial L}{\partial z_{ij}} \frac{\partial z_{ij}}{\partial W_{jk}} = \sum_{i=1}^N \frac{\partial L}{\partial z_{ij}} x_{ik}
\]

\[
\frac{\partial L}{\partial b_j} = \sum_{i=1}^N \frac{\partial L}{\partial z_{ij}} \frac{\partial z_{ij}}{\partial b_j} = \sum_{i=1}^N \frac{\partial L}{\partial z_{ij}}
\]

In matrix form:
\[
\frac{\partial L}{\partial \Wbf} = \left(\frac{\partial L}{\partial \Zbf}\right)^T \Xbf, \quad \frac{\partial L}{\partial \bbf} = \left(\frac{\partial L}{\partial \Zbf}\right)^T \one
\]
where $\Xbf$ is $N \times N_i$ with rows $\xbf_i^T$, and $\one$ is a vector of ones.

 \item
 Put all above together, describe how would you compute the gradient  $\partial L/\partial W_{jk}$ and $\partial L/\partial b_j$.

\textbf{Solution:}
Backpropagation steps:
\begin{enumerate}
\item Forward pass: Compute $\zbf_i$, $\ubf_i$, $\hat{y}_i$ for all $i$.
\item Compute $\frac{\partial L}{\partial \hat{y}_i} = -2(y_i - \hat{y}_i)$ for all $i$.
\item Compute $\frac{\partial L}{\partial u_{ij}} = \frac{\partial L}{\partial \hat{y}_i} \frac{a_j - \hat{y}_i}{S_i}$.
\item Compute $\frac{\partial L}{\partial z_{ij}} = \frac{\partial L}{\partial u_{ij}} u_{ij}(1 - u_{ij})$.
\item Compute $\frac{\partial L}{\partial W_{jk}} = \sum_{i=1}^N \frac{\partial L}{\partial z_{ij}} x_{ik}$ and $\frac{\partial L}{\partial b_j} = \sum_{i=1}^N \frac{\partial L}{\partial z_{ij}}$.
\end{enumerate}

\item Write a few lines of python code to implement the gradients  $\partial L/\partial \ubf$ (as in part (d)), given the gradient $\partial L/\partial \hat{\bf y}$.
Indicate how you represent the gradients.
Full credit requires that you avoid for-loops.

\textbf{Solution:}
\textbf{Representation:}
\begin{itemize}
\item $\partial L/\partial \hat{\ybf}$: 1D array of shape $(N,)$
\item $\partial L/\partial \Ubf$: 2D array of shape $(N, M)$
\item $\abf$: 1D array of shape $(M,)$
\item $\hat{\ybf}$: 1D array of shape $(N,)$
\item $\Ubf$: 2D array of shape $(N, M)$
\end{itemize}

\begin{python}
import numpy as np

def grad_L_wrt_u(dL_dyhat, a, yhat, U):
    """
    Compute gradient of L with respect to U.
    
    Args:
        dL_dyhat: 1D array of shape (N,) - dL/dyhat
        a: 1D array of shape (M,) - weights a_j
        yhat: 1D array of shape (N,) - predictions yhat_i
        U: 2D array of shape (N, M) - hidden activations u_{ij}
    
    Returns:
        dL_dU: 2D array of shape (N, M) - dL/dU
    """
    S = np.sum(U, axis=1)  # S_i = sum_j u_{ij}, shape (N,)
    dL_dU = (dL_dyhat[:, np.newaxis] / S[:, np.newaxis]) * (a[np.newaxis, :] - yhat[:, np.newaxis])
    return dL_dU
\end{python}

\end{enumerate}

\end{enumerate}
\end{document} 